\documentclass[10pt]{article}

\def\Type{LessonDJ}
\def\TypeNum{Workshop Week 13} %Lesson #
\def\TypeTitle{FTOC and Intro to Probability} %Lesson Title
\def\Semester{Fall 2021} %Not Used 
\def\Course{MATH 1190}%Course number and name
\def\Targets{    \target{I2}, \target{I3}, \target{F3} }

\usepackage{preambleDJ} 

%%%%%%%%%%%% BEGIN DOCUMENT %%%%%%%%%%%%%%%
%%%%%%%%%%%% BEGIN DOCUMENT %%%%%%%%%%%%%%%
%%%%%%%%%%%% BEGIN DOCUMENT %%%%%%%%%%%%%%%
%%%%%%%%%%%% BEGIN DOCUMENT %%%%%%%%%%%%%%%
%%%%%%%%%%%% BEGIN DOCUMENT %%%%%%%%%%%%%%%

\begin{document}
\pagestyle{otherpages}
\thispagestyle{firstpage}


\vs{.6}

\section{Warm-Up}
\begin{enumerate}
    \item Find the following anti-derivatives 
    \begin{enumerate}
        \item $\ds \int x^3 \;dx$\vfill
        \item $\ds \int \frac{1}{2x} \;dx$\vfill
        \item $\ds  \int 4\cdot3^x\;dx$\vfill
    \end{enumerate}
    \item Compute the following definite integral
    \[\int_2^5 \frac{x^4+1 + x3^x}{2x}dx\]\vfill\vfill
    \item Compute the following definite integrals:
    \begin{enumerate}
        \item $\ds \int_0^3 (2x+1)^2\,dx.$\vfill\vfill
         \item $\ds \int_1^4 \left(\frac 1x + 2\right)\left(3x + 1\right)\,dx.$\vfill\vfill
    \end{enumerate}    

\end{enumerate}

\newpage
 \section{More Definite Integrals}
   \begin{enumerate}
    \item Consider the following integral:
    \[\int_{-1}^3\frac{1}{x^4}dx\]
    Does the fundamental theorem of calculus apply? If not, explain why not. If so, find the definite integral.\\\vfill
    \item Consider the function $v(t) = 4t^3-16t^2+16t$ shown below:
    \begin{center}
        \begin{tikzpicture}
          \begin{axis}[
            xmin=-2.5,        % Minimum x-value
            xmax=2.5,         % Maximum x-value
            ymin=-8,         % Minimum y-value
            ymax=8,          % Maximum y-value
            axis lines=middle, % Draw the x and y axes in the middle
            x label style={anchor=west},
            xlabel={$t$},     % Label for the x-axis
            ylabel={$v(t)$},     % Label for the y-axis
            xtick       = {-2,-1,1,2},
            ytick       = {-8,-4,4,8},
            %yticklabels = {$2$,$4$,$6$,$8$},
            axis line style={thick, <->}, % Add arrows at the end of the axes
            width = 2in, height=2in,
            %restrict y to domain=-6.5:6.5
          ]
          %\addplot[domain=-2.5:2.5, samples=400, blue, ultra thick] {x^3};
              \addplot[name path=f1, domain=-2.5:2.5, blue, ultra thick, <->] {4*x^3-16*x^2+16*x};
%  \path[name path=axis] (axis cs:-2.5,0) -- (axis cs:2.5,0);
%     \addplot [
%        thick,
%        color=blue,
%        fill=blue, 
%        fill opacity=0.1
%    ]
%    fill between[
%        of=f1 and axis,
%        soft clip={domain=-1:2},
%   ];
          \end{axis}
        \end{tikzpicture}
    \end{center}

    \begin{enumerate}
        \item[a)] What is the area bounded by $v$ and the $t$-axis from $t = 1$ to $t = 2$?\vfill
        \item[b)] What is the net area bounded by $v$ and the $t$-axis from $t = -1$ to $t = 1$?\vfill
        \item[c)] What is the net area bounded by $v$ and the $t$-axis from $t = -1$ to $t = 2$?\vfill
    \end{enumerate}

    \item Suppose $v$ represented the velocity of a car as a function of time. What do the quantities you calculated in the previous part represent?\vfill
    
\end{enumerate}

\end{document}